\section{Explicit spectral integrals for distinct phase coefficients}
\label{sec:distinct-phase-integrals}

For distinct phase coefficients within each group, the complete integral
set can be chosen directly from spectral values and derivatives. This
provides an explicit specialization of \cref{thm:canonical-independence}.
The selection points are determined by
the fixed diagonal entries of the pencil. In the $(\nu,\lambda)$
plane, the diagonal factor $\lambda-X_{hh}-\nu D_{hh}$ vanishes
on the straight line $\lambda=X_{hh}+\nu D_{hh}$,
$1\leq h\leq d$. From \cref{eq:canonical-X,eq:canonical-pencil},
these are the three families
\begin{equation}\label{eq:canonical-diagonal-lines}
\begin{aligned}
 \lambda&=\nu-2b_\ell^L,&&1\leq\ell\leq L,\\
 \lambda&=-\nu-2b_m^M,&&1\leq m\leq M,\\
 \lambda&=b_n^N,&&1\leq n\leq N.
\end{aligned}
\end{equation}
Thus a family is simply a list of parallel straight lines: there are
$L$ lines of slope $1$, $M$ of slope $-1$, and $N$ horizontal lines.
The $\lambda$-intercept is the value at $\nu=0$, namely $X_{hh}$;
the three lists of intercepts are $-2b_\ell^L,-2b_m^M,b_n^N$.

An intersection $q=(\lambda_q,\nu_q)$ is fixed by these phase
coefficients. Evaluating the spectral determinant there gives the
function $F(q)=F(\lambda_q,\nu_q)$ of the evolving modes.
These values are conserved and commute by
\cref{prop:canonical-commutation}. Distinct intercepts within each
family ensure that at most one line from each family passes through
any point, so intersections involve either two or three lines.

\begin{theorem}\label{thm:distinct-phase-evaluations}
Suppose each of the three lists of phase coefficients
\[
 b_1^L,\ldots,b_L^L;\qquad
 b_1^M,\ldots,b_M^M;\qquad
 b_1^N,\ldots,b_N^N
\]
is pairwise distinct.
At every two-line intersection $q$, take $F(q)$; at every triple
intersection, take $F(q),\partial_\lambda F(q),\partial_\nu F(q)$.
These $p=LM+MN+NL$ functions are generically independent commuting
first integrals. Hence \cref{eq:canonical-H} is Liouville integrable
for every positive $L,M,N$. A complete local set can include $H$
wherever $\mathrm dH\ne0$.
\end{theorem}
\begin{proof}
There are $LM+MN+NL$ pairs of lines from different families.
Ordinarily each pair gives one evaluation. If three lines meet,
their three pairwise intersections coincide; the value and its two
first spectral derivatives provide three functions at that point.
Thus the total count remains $p$. The derivatives also commute,
by differentiating \cref{prop:canonical-commutation} with respect
to the fixed spectral arguments.

To prove independence, it suffices to find one nonzero Jacobian minor.
Use the $p$ upper off-diagonal mode entries $x_{hi}=X_{hi}$ as
column coordinates, with $1\leq h<i\leq d$ in different blocks;
use the proposed integrals as rows. Evaluate this Jacobian at the
test states $X_{hi}=0$, $X_{ih}=\epsilon$, keeping all phase
coefficients fixed. Here $\epsilon$ is a parameter selecting the
test state. We find the leading determinant as $\epsilon\to0$.

Write $l_h=\lambda-X_{hh}-\nu D_{hh}$, $1\leq h\leq d$,
and let $R_q$ be the product of the nonzero $l_h(q)$ over
$1\leq h\leq d$. At a two-line intersection, two diagonal
factors vanish. A determinant term linear in $x_{hi}$ can remove
both zeros by pairing it with its reverse entry $\epsilon$;
this gives $-\epsilon R_q$ only for the edge joining those two
lines. All other columns vanish to higher order. Dividing this
row by $\epsilon$ therefore leaves the single nonzero leading
entry $-R_q$.

At a triple intersection, three diagonal factors vanish. The Jacobian
row for $F(q)$ first detects a three-edge cycle, of order $\epsilon^2$.
A spectral derivative removes one vanishing factor, so
$F_\lambda(q)$ and $F_\nu(q)$ detect two-edge cycles of order
$\epsilon$. Order the three relevant columns as $L$--$M$,
$L$--$N$, $M$--$N$. After dividing these three rows by
$\epsilon^2,\epsilon,\epsilon$, respectively, determinant
differentiation gives the leading block
\begin{equation}\label{eq:canonical-confluent-block}
 R_q\begin{pmatrix}0&-1&0\\-1&-1&-1\\0&-1&1\end{pmatrix},
 \qquad \det=-R_q^3\ne0.
\end{equation}
Leading entries on edges outside the corresponding intersection
vanish, because a vanishing diagonal factor remains. Consequently
the full normalized Jacobian is block diagonal: one nonzero scalar
block for each two-line intersection and one invertible three-by-three
block for each triple. Its determinant is nonzero. The unscaled
minor is therefore a nonzero polynomial in $\epsilon$, proving
independence on a dense open set of the mode phase space. Since $H$
commutes with this complete set, it is locally a function of these
integrals and can replace a suitable member wherever $\mathrm dH\ne0$.
\end{proof}

The $d-1$ independent power charges can be retained in a complete set,
leaving
\[
 G=p-d+1=LM+MN+NL-L-M-N+1.
\]
spectral integrals. Away from triple intersections, there is a simple
group-symmetric selection: regard the modes as edges of the complete
tripartite graph, choose any spanning tree, and retain the evaluations
on its $G$ complementary edges. Indeed, at the proof's test state the
charge differentials form a weighted incidence matrix, whose tree
minor is nonzero; the complementary spectral rows give the remaining
nonzero diagonal block. At triple intersections, differential row
selection from the confluent set gives the required complement.
For $N=1$, the count becomes $L+M$ power charges and $LM$ spectral
integrals.

